problem:  
Let \((M_i, g_i)\) be a sequence of compact \(n\)-dimensional Riemannian manifolds with almost nonnegative sectional curvature, i.e.  
\[
\mathrm{sec}_{g_i} \ge -\varepsilon_i, \qquad \varepsilon_i \to 0,
\]
and assume \(\mathrm{diam}(M_i) \le 1\). Suppose that \((M_i, g_i)\) converge in the Gromov–Hausdorff sense to a compact length space \(X\).  
By the results of Cheeger–Fukaya–Gromov, if curvature has two-sided bounds, collapsing regions carry local nilpotent structures, while Yamaguchi’s fibration theorem implies that under lower curvature bounds such collapse yields local Riemannian submersion models with Alexandrov base spaces of curvature \(\ge 0\). However, the global structure of the limit space \(X\) is still poorly understood.

**Problem (Global resolution question for almost nonnegative collapse):**  
Does every such Gromov–Hausdorff limit space \(X\) admit a realization, up to isometry, as the metric quotient  
\[
X \cong N / G,
\]
where \(N\) is a smooth compact manifold carrying a Riemannian metric with nonnegative sectional curvature and \(G\) is a compact Lie group acting isometrically, effectively, and locally freely on \(N\), such that the *length metric* on \(X\) coincides with the induced quotient metric on \(N/G\)?  

In other words, can every Alexandrov space that arises as a Gromov–Hausdorff limit of almost nonnegatively curved manifolds be resolved as the quotient of a smooth nonnegatively curved Riemannian manifold by a compact Lie group action?

Why is it a "good" problem:  
This question probes the global geometric and topological structure of limits in collapsing sequences with lower sectional curvature bounds, a domain where local models (nilpotent or submersion-type structures) are known but global smooth descriptions are lacking. It aims to test whether the singular Alexandrov limits of smooth almost nonnegatively curved manifolds can always be represented as smooth quotients with genuine nonnegative curvature. A positive answer would provide a powerful unification of Alexandrov geometry and classical Riemannian geometry via Lie group actions; a negative answer would reveal fundamentally new singular geometries unreachable by smooth curvature models.  

The problem is conceptually simple yet requires bridging deep areas—collapsing theory, transformation group theory, topology of curvature, Alexandrov geometry, and metric structure theory—making it a profound and far-reaching research question that typifies what a “good” problem should be.